\section{From Sources to Temporal Evidence}
\label{sec:sources}

An occurrence date says when a state began.
A duration gives separate start and end dates.
Table~\ref{tab:date-types} shows how we represent both.

\begin{table}[b]
\centering\small
\begin{tabular}{@{}p{.28\columnwidth}p{.65\columnwidth}@{}}
\toprule
Source meaning & Representation\\
\midrule
Started in September & An unknown start between that month's first and last dates.\\[3pt]
Served from 2001 to 2005 & Separate uncertain start and end dates.\\
\bottomrule
\end{tabular}
\caption{Date precision and state duration carry different information. Only an occurrence range bounds one unknown event date.}
\label{tab:date-types}
\end{table}

\paragraph{Occurrence precision.}
A day gives a point on the time axis.
A month or year bounds an occurrence within that calendar period.
Expressions such as \emph{around 1713} require additional evidence before they supply finite bounds.
We preserve the literal date expression and its source span.
The query adapter applies the same convention to the visible question.
A question about a whole year therefore remains a period query.
Our predicate asks whether a claim applies \emph{somewhere} within that period.

\paragraph{Replacement and coexistence.}
A directed replacement decision connects values of the same subject and role.
For office succession, the organization and office define the slot; people supply its values.
For affiliation, a new organization can coexist with an older affiliation.
The Chicago Stadium passage in TimeQA, for example, gives overlapping tenancies for two teams~\citep{chen2021timeqa}.
The source adapter records explicit succession separately from ordinary appointments or joins.

\paragraph{Explicit ends.}
A stated departure supplies an end event targeting its own claim.
The compiler treats this hidden event as a direct replacement.
Every permitted end must follow every permitted start: its lower bound exceeds the start's upper bound.
This reduction preserves gaps between tenures.
It also handles a departure without a named successor.
Appendix~\ref{app:lifecycle} proves equivalence to the corresponding half-open lifetime.

\paragraph{Evidence units.}
A paragraph can contain several states, so temporal selection operates on source excerpts.
Each modeled claim preserves exact character offsets and its supporting sentence for audit.
Uncovered text remains available through a shared fallback policy.
A loan's planned expiry never becomes an observed departure.
An incomplete temporal parse retains the complete source text.
Readers receive original excerpts, date bounds, and support tags.
Source-derived subject and relation labels preserve context for short succession clauses.

\paragraph{Support tags.}
Modeled claims are guaranteed, unresolved, or unsupported for the requested period.
Unresolved claims have possible support without guaranteed support.
Unmodeled excerpts carry a separate \emph{unknown} tag.
Their fixed fallback eligibility lets us isolate uncertainty in modeled dates.

\paragraph{A source-grounded example.}

TimeQA's Warburg Institute passage names three successive directors~\citep{chen2021timeqa}.
Henri Frankfort started in 1949, Gertrud Bing in 1955, and Ernst Gombrich in 1959.
A deterministic clause adapter recovers these three starts and their two direct succession links from the paragraph.
The passage separately dates Bing's earlier organizational membership to 1922.
That membership date does not start her directorship.

Consider the illustrative period January 1954 through June 1955.
Frankfort has guaranteed support somewhere within that period.
Bing has possible support because her appointment might precede July.
Gombrich has no support within the period.
Rank Frankfort first and Bing second.
A one-item context always contains Frankfort.
A two-item context changes when Bing's start moves from January to December 1955.
Both assignments satisfy the source's year-level precision.
The same possible pool thus yields different stability decisions at different budgets.
The source is natural; the query, ranking, and assignments illustrate the theorem.
